% ============================================================================
% WORKBOOK — Additional: Financial Literacy — Budgeting, Saving, and Investing
% State-specific (TX)
% Practice-focused reinforcement for the study guide topic
% ============================================================================
\section{Financial Literacy --- Budgeting, Saving, and Investing}
\topicTitle{Financial Literacy --- Budgeting, Saving, and Investing}

\begin{quickReview}{Budgeting Basics}
A \textbf{budget} is a plan for your money: \textbf{Income} $-$ \textbf{Expenses} $=$ \textbf{Savings}.

\begin{itemize}[leftmargin=*]
    \item \textbf{Fixed expenses} stay the same each month (rent, phone plan).
    \item \textbf{Variable expenses} change month to month (food, entertainment).
    \item \textbf{Savings rate} $= \dfrac{\text{savings}}{\text{income}} \times 100\%$.
    \item \textbf{Investing} means putting money into assets (stocks, bonds) hoping it grows; it carries \textbf{risk}.
\end{itemize}

\medskip
\textbf{Example:} Income $= \$1{,}200$. Fixed $= \$700$. Variable $= \$350$. Savings $= 1{,}200 - 700 - 350 = \$150$. Rate $= \frac{150}{1{,}200} = 12.5\%$.
\end{quickReview}

% ── Warm-Up ──────────────────────────────────────────────────────────────────
\begin{practiceBox}[funGreen]{\faSun~Warm-Up}

\practiceHeader[funGreen]{How Much Is Left to Save?}

\begin{multicols}{2}
\prob Income: $\$1{,}000$. Expenses: $\$800$. Savings $=$ \answerBlank[2cm]
\answerExplain{$\$200$}{Subtract expenses from income: $1{,}000 - 800 = \$200$.}

\prob Income: $\$2{,}500$. Expenses: $\$2{,}100$. Savings $=$ \answerBlank[2cm]
\answerExplain{$\$400$}{$2{,}500 - 2{,}100 = \$400$.}

\prob Income: $\$1{,}800$. Fixed: $\$900$. Variable: $\$600$. Savings $=$ \answerBlank[2cm]
\answerExplain{$\$300$}{$1{,}800 - 900 - 600 = \$300$.}

\prob Income: $\$3{,}000$. Fixed: $\$1{,}500$. Variable: $\$1{,}200$. Savings $=$ \answerBlank[2cm]
\answerExplain{$\$300$}{$3{,}000 - 1{,}500 - 1{,}200 = \$300$.}

\prob Income: $\$2{,}000$. Total expenses: $\$1{,}750$. Savings rate $=$ \answerBlank[2cm]
\answerExplain{$12.5\%$}{Savings $= 2{,}000 - 1{,}750 = \$250$. Rate $= \frac{250}{2{,}000} = 0.125 = 12.5\%$.}

\prob Income: $\$4{,}000$. Savings: $\$600$. Savings rate $=$ \answerBlank[2cm]
\answerExplain{$15\%$}{Rate $= \frac{600}{4{,}000} = 0.15 = 15\%$.}
\end{multicols}
\end{practiceBox}

% ── Practice A: Building a Budget ────────────────────────────────────────────
\begin{practiceBox}{Build a Monthly Budget}

\prob Jake earns $\$2{,}400$/month. His rent is $\$850$, phone is $\$50$, groceries $\$350$, and entertainment $\$200$. Which expenses are fixed and which are variable? \answerBlank[5cm]
\answerExplain{Fixed: rent, phone; Variable: groceries, entertainment}{Rent and phone cost the same every month (fixed). Groceries and entertainment amounts can change (variable).}

\prob Using Jake's budget from Problem 7, how much can he save per month? \answerBlank[3cm]
\answerExplain{$\$950$}{Total expenses: $850 + 50 + 350 + 200 = \$1{,}450$. Savings: $2{,}400 - 1{,}450 = \$950$.}

\prob What is Jake's savings rate as a percent? Round to the nearest whole percent. \answerBlank[3cm]
\answerExplain{$40\%$}{$\frac{950}{2{,}400} \approx 0.396 \approx 40\%$.}

\prob Maria's monthly income is $\$1{,}600$. She wants to save $20\%$. How much can she spend on expenses? \answerBlank[3cm]
\answerExplain{$\$1{,}280$}{Savings: $1{,}600 \times 0.20 = \$320$. Expenses: $1{,}600 - 320 = \$1{,}280$.}

\prob Your food budget is $\$400$/month. One month you spend $\$450$. By what percent did you go over budget? \answerBlank[3cm]
\answerExplain{$12.5\%$}{Over by $\$50$. Percent: $\frac{50}{400} = 0.125 = 12.5\%$.}
\end{practiceBox}

% ── Practice B: Savings Goals ────────────────────────────────────────────────
\begin{practiceBox}[funTeal]{Reaching a Savings Goal}

\prob You want to save $\$1{,}200$ in $6$ months. How much per month? \answerBlank[3cm]
\answerExplain{$\$200$/month}{Divide the goal by the number of months: $1{,}200 \div 6 = \$200$ per month.}

\prob You want to save $\$3{,}000$ in $1$ year. How much per week (assume $52$ weeks)? Round to the nearest cent. \answerBlank[3cm]
\answerExplain{$\$57.69$/week}{$3{,}000 \div 52 \approx \$57.69$ per week.}

\prob You save $\$75$/month. In how many months will you reach $\$900$? \answerBlank[3cm]
\answerExplain{$12$ months}{$900 \div 75 = 12$ months.}

\prob You save $\$200$/month at $3\%$ simple annual interest. After $12$ months you've deposited $\$2{,}400$. About how much interest will you earn over the year? \answerBlank[3cm]
\answerExplain{About $\$36$}{On average your balance is about $\$1{,}200$ over the year. $I \approx 1{,}200 \times 0.03 \times 1 = \$36$.}

\prob You have $\$500$ and add $\$100$/month. How many months until you reach $\$1{,}500$? \answerBlank[3cm]
\answerExplain{$10$ months}{You need $1{,}500 - 500 = \$1{,}000$ more. $1{,}000 \div 100 = 10$ months.}
\end{practiceBox}

% ── Practice C: Saving vs. Investing ─────────────────────────────────────────
\begin{practiceBox}[funOrange]{Saving vs.\ Investing}

\prob A savings account earns $2\%$ per year. How much interest does $\$5{,}000$ earn in $1$ year? \answerBlank[3cm]
\answerExplain{$\$100$}{$I = 5{,}000 \times 0.02 \times 1 = \$100$.}

\prob An investment averages $7\%$ growth per year. How much does $\$5{,}000$ grow in $1$ year? \answerBlank[3cm]
\answerExplain{$\$350$}{$5{,}000 \times 0.07 = \$350$.}

\prob How much more does the investment earn compared to the savings account? \answerBlank[3cm]
\answerExplain{$\$250$ more}{$350 - 100 = \$250$ more. However, the investment carries risk of losing money.}

\prob Name one reason to choose a savings account over an investment. \answerBlank[5cm]
\answerExplain{Safety/low risk}{Savings accounts are low-risk --- your money is protected. Investments can lose value.}
\end{practiceBox}

% ── Word Problem ─────────────────────────────────────────────────────────────
\begin{practiceBox}[funBlue]{\faGlobe~Word Problems}

\wordProblem{Alicia earns $\$2{,}800$/month. She spends $\$1{,}200$ on rent, $\$400$ on food, $\$150$ on transportation, and $\$250$ on other expenses. What percent of her income does she save? Round to the nearest whole percent.}{}
\answerExplain{$29\%$}{Total expenses: $1{,}200 + 400 + 150 + 250 = \$2{,}000$. Savings: $2{,}800 - 2{,}000 = \$800$. Rate: $\frac{800}{2{,}800} \approx 0.286 \approx 29\%$.}

\wordProblem{Ben puts $\$3{,}000$ in a savings account at $2.5\%$ simple interest. How much will he have after $4$ years?}{}
\answerExplain{$\$3{,}300$}{$I = 3{,}000 \times 0.025 \times 4 = \$300$. Total $= 3{,}000 + 300 = \$3{,}300$.}
\end{practiceBox}

% ── Challenge ────────────────────────────────────────────────────────────────
\begin{challengeBox}
\prob A family earns $\$5{,}200$/month. They budget $30\%$ for housing, $15\%$ for food, $10\%$ for transportation, $5\%$ for entertainment, and save the rest. How much do they save, and what is their savings rate? \answerBlank[3cm]
\answerExplain{$\$2{,}080$ saved, $40\%$ rate}{Used: $30 + 15 + 10 + 5 = 60\%$. Savings: $100\% - 60\% = 40\%$. Amount: $5{,}200 \times 0.40 = \$2{,}080$.}
\end{challengeBox}

\encouragement{Great job with budgeting! Planning your money wisely is one of the most important life skills!}
